% Source: https://artofproblemsolving.com/community/c3h1936678_two_squares
% Stable identity: aops-topic-1936678-q01
\begin{asy}
usepackage("amsmath");
size(5cm);
defaultpen(fontsize(10pt));
real x=8, y=15;
draw((0,0)--(x,0)--(x,x)--(0,x)--cycle,linewidth(0.8));
draw((11,0)--(11+y,0)--(11+y,y)--(11,y)--cycle,linewidth(0.8));
label("$A$",(x/2,x/2)); label("$B$",(11+y/2,y/2));
label("$x$",(x/2,0),S); label("$y$",(11+y/2,0),S);
label("\mbox{}",(min(currentpicture,user=true)+max(currentpicture,user=true))/2,black); // keeps the text after the figure black when LaTeX typesets the labels inline
\end{asy}

The perimeter condition gives $4y-4x=28$, so $y-x=7$. The area condition gives $y^2-x^2=161$. Factoring the difference of squares,
\[(y-x)(y+x)=161\quad\Longrightarrow\quad 7(x+y)=161\quad\Longrightarrow\quad x+y=\boxed{23}.\]
(Indeed $x=8$ and $y=15$: perimeters $32$ and $60$, areas $64$ and $225$.)
